\documentclass{beamer}

\mode<presentation>
\title{SoC1\&2: RISC-V Post-Route Optimization with Innovus}
\author[AS]{Andreas Siggelkow \newline \texttt{andreas.siggelkow@rwu.de}}
\date[2026]{\today}
\subtitle{RWU}
\institute{Elektrotechnik und Informatik}

%%%%%%%%%%%%%%%%%%%%%%%%%%%%%%%%
% Logo
%%%%%%%%%%%%%%%%%%%%%%%%%%%%%%%%
\pgfdeclareimage[height=1cm]{hslogo}{./hs}
\logo{\pgfuseimage{hslogo}}
%%%%%%%%%%%%%%%%%%%%%%%%%%%%%%%%
% Macros
%%%%%%%%%%%%%%%%%%%%%%%%%%%%%%%%
\usepackage[T1]{fontenc}
\usepackage{lmodern}
\usepackage{textcomp}
\definecolor{links}{HTML}{2A1B81}
\hypersetup{colorlinks,linkcolor=,urlcolor=links}
\usepackage{graphicx}
\usepackage{ragged2e, tabularx}
\newcolumntype{Y}{>{\RaggedRight\arraybackslash}X}
\usetheme{Singapore}
\usecolortheme{seahorse}
\usefonttheme[onlylarge]{structuresmallcapsserif}
\usefonttheme[onlysmall]{structurebold}
\setbeamertemplate{footline}[frame number]
\usepackage{listings}
\newcommand{\astitle}{nix}
\newcommand{\assubtitle}{nix}
\newcommand{\asfolie}[3]{\begin{frame}\frametitle{#1}\begin{block}{#2}{#3}\end{block}\end{frame}}

% Commands and tool output are set with listings: dashes, $, _, \ and quotes
% appear exactly as they have to be typed.
\definecolor{lstbg}{HTML}{F2F2F7}
\definecolor{lstcomment}{HTML}{5A5A8A}
\lstset{
  basicstyle=\ttfamily\scriptsize,
  columns=fullflexible,
  keepspaces=true,
  upquote=true,
  breaklines=true,
  breakatwhitespace=false,
  backgroundcolor=\color{lstbg},
  frame=none,
  xleftmargin=2pt,
  xrightmargin=2pt,
  aboveskip=4pt,
  belowskip=4pt,
  showstringspaces=false
}
\lstdefinestyle{sh}{
  morecomment=[l]{\#},
  commentstyle=\color{lstcomment}\itshape
}
\lstdefinestyle{out}{basicstyle=\ttfamily\tiny}
\newcommand{\cmd}[1]{\texttt{#1}}

\begin{document}
%%%%%%%%%%%%%%%%%%%%%%%%%%%%%%%%
% First pages
%%%%%%%%%%%%%%%%%%%%%%%%%%%%%%%%
\setcounter{tocdepth}{1}

\begin{frame}
  \titlepage
\end{frame}

\begin{frame}
  \frametitle{Outline}
  \tableofcontents
\end{frame}

%%%%%%%%%%%%%%%%%%%%%%%%%%%%%%%%
% Overview
%%%%%%%%%%%%%%%%%%%%%%%%%%%%%%%%
\section{Overview}
\renewcommand{\astitle}{Overview}
\begin{frame}
\frametitle{\astitle}
\textcolor{red}{Overview}

\url{https://github.com/asiggel/rwu-rv64iV2.0}
\end{frame}

\renewcommand{\assubtitle}{Where in the Flow?}
\begin{frame}
\frametitle{\astitle}
\begin{block}{\assubtitle}
  \footnotesize
  \begin{tabularx}{\linewidth}{@{}llY@{}}
    \textbf{step} & \textbf{tool / script} & \textbf{result}\\\hline
    01 RTL simulation   & xrun, \cmd{scripts/xcelium/}  & functionally correct RTL\\
    02 synthesis + DFT  & Genus, \cmd{scripts/genus/}   & netlist, scan chain\\
    03 place \& route   & Innovus, \cmd{innovus\_pnr.sh} & routed layout\\
    04 gate-level sim   & xrun, \cmd{scripts/gatesim/}  & netlist ok\\
    \textbf{05 post-route} & \textbf{Innovus}, \cmd{innovus\_post.sh} & final layout, hand-off, GDS\\
    06 extraction + STA & Quantus, Tempus               & sign-off timing\\
    07 LEC              & Conformal                     & RTL $\equiv$ layout netlist\\
    08 DRC/LVS/ERC      & Pegasus/Calibre               & geometry ok\\
  \end{tabularx}
  \vspace{4pt}
  \begin{itemize}
    \item Numbers = HowTo folders, in the order they run
    \item Innovus twice: 03 builds the layout up to the routing, 05
      finishes it
  \end{itemize}
\end{block}
\end{frame}

\renewcommand{\assubtitle}{Why Optimize after Routing?}
\begin{frame}
\frametitle{\astitle}
\begin{block}{\assubtitle}
  \begin{itemize}
    \item Before routing the timing uses estimated wires; after routing
      the wires are real: longer detours, coupling to the neighbours
    \item The clock tree (03, CTS) adds skew: a flip-flop to flip-flop
      path with almost no logic can now violate \textbf{hold}
    \item Hold violations do not depend on the clock frequency: a chip
      with a hold violation does not work at any speed
    \item Post-route optimization repairs what routing broke:
    \begin{itemize}
      \item setup: resize cells, insert buffers on long wires
      \item hold: insert delay cells in too-fast data paths
      \item DRV: max transition / capacitance / fanout
    \end{itemize}
    \item The router makes the small ECO routes itself
  \end{itemize}
\end{block}
\end{frame}

\renewcommand{\assubtitle}{Why a Separate Step?}
\begin{frame}
\frametitle{\astitle}
\begin{block}{\assubtitle}
  \begin{itemize}
    \item The optimization inserts and resizes cells. Everything that
      fills the layout must come \emph{after} it:
    \begin{itemize}
      \item core fillers (in the gaps that are left)
      \item metal fill (in the empty metal that is left)
    \end{itemize}
    \item The hand-off files (netlists, DEF, GDS) must show the
      optimized layout: Quantus, Tempus, LEC, GLS and DRC/LVS all check
      the \emph{same} layout
    \item The P\&R run (03) takes long; the post-route step starts from
      its checkpoint \cmd{cc\_chip\_top\_routed.db} and can be repeated
      alone
  \end{itemize}
\end{block}
\end{frame}

\renewcommand{\assubtitle}{Inputs and Outputs}
\begin{frame}
\frametitle{\astitle}
\begin{block}{\assubtitle}
  \scriptsize
  \begin{tabularx}{\linewidth}{@{}lY@{}}
    \textbf{In}                       & \\\hline
    \cmd{build\_innovus/cc\_chip\_top\_routed.db} & routed layout from 03\\[1pt]
    \cmd{as\_top\_mem\_scan.sdc}      & constraint mode scan (for the hold view)\\[1pt]
    \cmd{\$XFAB\_PATHS}               & \cmd{CORE\_FILLERS}, \cmd{FILL\_LAYERS},
                                        GDS variables\\[4pt]
    \textbf{Out} (same directory)     & \\\hline
    \cmd{cc\_chip\_top\_postopt.db}   & checkpoint after the optimization\\[1pt]
    \cmd{cc\_chip\_top\_final.db}     & final layout (fillers, metal fill)\\[1pt]
    \cmd{*\_post\_\{setup,hold\}\_pre.rpt} & worst paths before\\[1pt]
    \cmd{*\_final\_timing.rpt}, \cmd{*\_final\_hold.rpt} & worst paths after\\[1pt]
    \cmd{timing\_pre/}, \cmd{timing\_post/} & \cmd{time\_design} reports per view\\[1pt]
    \cmd{*\_drc/antenna/density/connectivity.rpt} & physical checks\\[1pt]
    \cmd{\_routed.v}, \cmd{\_routed.def} & hand-off to Quantus, Tempus, GLS, LEC\\[1pt]
    \cmd{\_lvs.v}, \cmd{cc\_chip\_top.gds} & hand-off to DRC/LVS\\
  \end{tabularx}
\end{block}
\end{frame}

%%%%%%%%%%%%%%%%%%%%%%%%%%%%%%%%
% Run
%%%%%%%%%%%%%%%%%%%%%%%%%%%%%%%%
\section{Run}
\renewcommand{\astitle}{Run}
\begin{frame}
\frametitle{\astitle}
\textcolor{red}{Run the post-route step}
\end{frame}

\renewcommand{\assubtitle}{Batch}
\begin{frame}[fragile]
\frametitle{\astitle}
\begin{block}{\assubtitle}
  \footnotesize
\begin{lstlisting}[style=sh,basicstyle=\ttfamily\small]
cd RWU_RV/RV_NoPipeline
./scripts/innovus/innovus_pnr.sh -b      # 03, once
./scripts/innovus/innovus_post.sh -b     # 05
\end{lstlisting}
% VERIFY-LAB: replace with the output of a real run
\begin{lstlisting}[style=out]
=== [innovus-post] 1/3 run 02d_rv64_postroute.tcl up to 'gds' (batch, extraction medium)
=== [innovus-post] 2/3 checkpoints
  opt: cc_chip_top_postopt.db
  finish: cc_chip_top_final.db
=== [innovus-post] 3/3 result
  worst setup slack: before ..., after ...
  worst hold  slack: before ..., after ...
  drc:              ...
  antenna:          ...
  density:          ...
  connectivity:     ...
  hand-off:         cc_chip_top_routed.v / _routed.def / _lvs.v
  GDS:              cc_chip_top.gds (...)
  GDS merge:        cells used: ..., missing: 0
PASS  post-route step of cc_chip_top up to 'gds'
\end{lstlisting}
\end{block}
\end{frame}

\renewcommand{\assubtitle}{Options}
\begin{frame}
\frametitle{\astitle}
\begin{block}{\assubtitle}
  \footnotesize
  \begin{tabularx}{\linewidth}{@{}lY@{}}
    \cmd{-b}, \cmd{-{}-batch}        & no GUI (default: Innovus GUI)\\[2pt]
    \cmd{-s}, \cmd{-{}-stop STAGE}   & stop after \cmd{opt}, \cmd{finish} or
                                       \cmd{gds} (default)\\[2pt]
    \cmd{-x}, \cmd{-{}-extract E}    & extraction effort for the optimization:
                                       \cmd{medium} (default), \cmd{high},
                                       \cmd{signoff} (Quantus engine, licence)\\[2pt]
    \cmd{-o}, \cmd{-{}-open STAGE}   & run nothing, open the checkpoint\\[2pt]
    \cmd{-h}, \cmd{-{}-help}         & help\\
  \end{tabularx}
  \vspace{6pt}
  \begin{itemize}
    \item Every run starts again from \cmd{cc\_chip\_top\_routed.db}; the
      P\&R checkpoints of 03 stay
    \item \cmd{-b -s opt}: only the timing repair, compare before/after
    \item A new P\&R run (03) deletes the results of this step
  \end{itemize}
\end{block}
\end{frame}

\renewcommand{\assubtitle}{Reading the Result}
\begin{frame}
\frametitle{\astitle}
\begin{block}{\assubtitle}
  \footnotesize
  \begin{tabularx}{\linewidth}{@{}lY@{}}
    hold after  & \textbf{must} be $\geq 0$ in every hold view, also
                  \cmd{av\_scan\_fast} (scan shift)\\[2pt]
    setup after & $\geq 0$ in \cmd{av\_slow}, \cmd{av\_jtag\_slow}. Hold
                  buffers cost a little setup: compare with before\\[2pt]
    DRV         & \cmd{timing\_post/}: max transition / capacitance left?\\[2pt]
    antenna     & near 0 (diodes during routing, 03)\\[2pt]
    density     & minimum density reached on every layer\\[2pt]
    GDS merge   & \cmd{missing: 0}\\
  \end{tabularx}
  \vspace{4pt}
  \begin{itemize}
    \item FAIL = Innovus error, negative slack after the optimization or
      a missing output (exit code 1)
    \item The numbers here use Innovus' fast extraction; the sign-off
      numbers come from Quantus + Tempus (06)
  \end{itemize}
\end{block}
\end{frame}

%%%%%%%%%%%%%%%%%%%%%%%%%%%%%%%%
% Under the hood
%%%%%%%%%%%%%%%%%%%%%%%%%%%%%%%%
\section{Under the Hood}
\renewcommand{\astitle}{Under the Hood}
\begin{frame}
\frametitle{\astitle}
\textcolor{red}{\cmd{02d\_rv64\_postroute.tcl} step by step}
\end{frame}

\renewcommand{\assubtitle}{a) Views: Hold in Scan Shift}
\begin{frame}[fragile]
\frametitle{\astitle}
\begin{block}{\assubtitle}
\begin{lstlisting}[style=sh]
read_db cc_chip_top_routed.db
create_constraint_mode -name scan \
  -sdc_files [list as_top_mem.sdc as_top_mem_scan.sdc]
create_analysis_view -name av_scan_fast -constraint_mode scan \
  -delay_corner dc_fast
set_analysis_view -setup {av_slow av_typ av_jtag_slow} \
                  -hold  {av_fast av_jtag_fast av_scan_fast}
\end{lstlisting}
  \begin{itemize}
    \item Genus and P\&R know only \cmd{func} and \cmd{jtag\_load}
    \item Scan shift: one flip-flop drives the next directly, no logic in
      between -- the classic hold problem, must be fixed in the layout
    \item Setup in scan mode stays out: it only limits TCK in SCANTEST
      (06\_DokuSTA)
  \end{itemize}
\end{block}
\end{frame}

\renewcommand{\assubtitle}{b) Analysis Settings}
\begin{frame}[fragile]
\frametitle{\astitle}
\begin{block}{\assubtitle}
\begin{lstlisting}[style=sh]
set_db extract_rc_engine       post_route
set_db extract_rc_effort_level $EXTRACT_EFFORT   ;# medium | high | signoff
set_db extract_rc_coupled      true
set_db timing_analysis_type    ocv
set_db timing_analysis_cppr    both
set_db delaycal_enable_si      true
\end{lstlisting}
  \begin{itemize}
    \item The same analysis as the sign-off STA (06): OCV, CPPR, SI
    \item Otherwise Innovus repairs against different numbers than
      Tempus checks
    \item \cmd{signoff} = Quantus inside Innovus: closest to the sign-off,
      slower, needs a Quantus licence
  \end{itemize}
\end{block}
\end{frame}

\renewcommand{\assubtitle}{c) Optimization}
\begin{frame}[fragile]
\frametitle{\astitle}
\begin{block}{\assubtitle}
\begin{lstlisting}[style=sh]
time_design -post_route        -report_dir timing_pre ...
time_design -post_route -hold  -report_dir timing_pre ...
opt_design  -post_route -setup -hold
opt_design  -post_route -drv
time_design -post_route        -report_dir timing_post ...
time_design -post_route -hold  -report_dir timing_post ...
write_db cc_chip_top_postopt.db
\end{lstlisting}
  \begin{itemize}
    \item Before/after: the reports show what the optimization changed
    \item Setup and hold in one call: Innovus balances them (a hold
      buffer must not create a setup violation)
    \item DRV afterwards: transition and capacitance limits of the
      Liberty cells
  \end{itemize}
\end{block}
\end{frame}

\renewcommand{\assubtitle}{d) Core Fillers}
\begin{frame}[fragile]
\frametitle{\astitle}
\begin{block}{\assubtitle}
\begin{lstlisting}[style=sh]
add_fillers -base_cells $CORE_FILLERS     ;# after the last cell change
\end{lstlisting}
  \begin{itemize}
    \item Fillers close the gaps in the rows: wells, implants and the
      supply rails become continuous -- otherwise DRC errors
    \item Physical only: not in the STA/GLS netlist, but in the LVS
      netlist
    \item After the optimization: a filler would block the space a hold
      buffer needs
  \end{itemize}
\end{block}
\end{frame}

\renewcommand{\assubtitle}{e) Metal Fill}
\begin{frame}[fragile]
\frametitle{\astitle}
\begin{block}{\assubtitle}
\begin{lstlisting}[style=sh]
add_metal_fill -layers $FILL_LAYERS
check_metal_density -report cc_chip_top_density.rpt
\end{lstlisting}
  \begin{itemize}
    \item Chemical-mechanical polishing (CMP) needs an even metal density:
      every layer needs a minimum (and maximum) share of metal per window
    \item Fill = unconnected dummy metal in the empty areas
    \item Before the export: Quantus sees the fill (extra capacitance), the
      GDS carries it
    \item Some foundries want their own fill deck (Pegasus/Calibre) instead:
      VERIFY-LAB, then this step goes
  \end{itemize}
\end{block}
\end{frame}

\renewcommand{\assubtitle}{f) Hand-off (\texttt{02b\_rv64\_export.tcl})}
\begin{frame}[fragile]
\frametitle{\astitle}
\begin{block}{\assubtitle}
\begin{lstlisting}[style=sh]
write_db cc_chip_top_final.db
write_netlist cc_chip_top_routed.v                          ;# STA, GLS, LEC
write_def -floorplan -netlist -routing cc_chip_top_routed.def ;# Quantus
write_netlist -include_pg_ports -include_phys_insts cc_chip_top_lvs.v
\end{lstlisting}
  \begin{itemize}
    \item The names stay \cmd{\_routed.*}: the later steps read them
      unchanged
    \item Two netlists: for timing only the logic; for LVS also the supply
      ports and all physical-only cells, because LVS compares device by
      device with the layout
  \end{itemize}
\end{block}
\end{frame}

\renewcommand{\assubtitle}{g) Merged GDS (\texttt{02c\_rv64\_gds.tcl})}
\begin{frame}[fragile]
\frametitle{\astitle}
\begin{block}{\assubtitle}
\begin{lstlisting}[style=sh]
write_stream cc_chip_top.gds -map_file $STREAM_MAP -merge $GDS_FILES \
             -unit $GDS_UNIT -mode ALL -structure_name cc_chip_top
\end{lstlisting}
  \begin{itemize}
    \item Innovus writes wires, vias and placements (it knows the cells only
      as LEF); \cmd{-merge} adds the full cell layouts from the X-FAB GDS
    \item \cmd{-map\_file}: LEF layer $\rightarrow$ GDS layer (PDK)
    \item \cmd{cc\_chip\_top\_gds.rpt}: cells used in the design but found
      in none of the GDS files -- must be 0, else the GDS has empty boxes
    \item Before: one text label per pad (port name or VDD/GND/VDDO/GNDO)
      on \cmd{LABEL\_LAYER}, for LVS
  \end{itemize}
\end{block}
\end{frame}

\renewcommand{\assubtitle}{h) Physical Checks on the Final Layout}
\begin{frame}[fragile]
\frametitle{\astitle}
\begin{block}{\assubtitle}
  \footnotesize
  Added by \cmd{innovus\_post.sh} (stages \cmd{finish}, \cmd{gds}):
\begin{lstlisting}[style=sh]
verify_connectivity -type all > cc_chip_top_connectivity.rpt   # in the flow
verify_drc -limit 100000 -report cc_chip_top_drc.rpt
verifyProcessAntenna -report cc_chip_top_antenna.rpt
\end{lstlisting}
  \begin{itemize}
    \item Same checks as after routing (03), now with hold buffers,
      fillers and metal fill
    \item Innovus' own checks with the LEF rules; sign-off is the X-FAB
      rule deck (08\_DokuPV)
  \end{itemize}
\end{block}
\end{frame}

%%%%%%%%%%%%%%%%%%%%%%%%%%%%%%%%
% Pitfalls
%%%%%%%%%%%%%%%%%%%%%%%%%%%%%%%%
\section{Pitfalls}
\renewcommand{\astitle}{Pitfalls}
\begin{frame}
\frametitle{\astitle}
\textcolor{red}{Pitfalls}
\end{frame}

\renewcommand{\assubtitle}{No Routed Checkpoint}
\begin{frame}[fragile]
\frametitle{\astitle}
\begin{block}{\assubtitle}
\begin{lstlisting}
ERROR: .../build_innovus/cc_chip_top_routed.db not found
       run ./scripts/innovus/innovus_pnr.sh -b first
\end{lstlisting}
  \begin{itemize}
    \item 03 has not run, or only up to an earlier stage (\cmd{-s fp})
    \item A P\&R run deletes the old post-route results: run 05 again
      after every run of 03
  \end{itemize}
\end{block}
\end{frame}

\renewcommand{\assubtitle}{Hold Not Fixed}
\begin{frame}
\frametitle{\astitle}
\begin{block}{\assubtitle}
  \begin{itemize}
    \item Hold still negative after \cmd{opt\_design}: look at the path in
      \cmd{cc\_chip\_top\_final\_hold.rpt}
    \begin{itemize}
      \item in \cmd{av\_scan\_fast}: scan path from one flip-flop to the
        next; large clock skew between them? (CTS, scan reorder, DFT-01)
      \item at a macro or a pad: the SDC may be wrong (input/output delay)
      \item no space for the delay cells: the core is too full
    \end{itemize}
    \item Setup much worse after the hold fixing: hold and setup on the
      same path conflict -- usually skew, not logic
  \end{itemize}
\end{block}
\end{frame}

\renewcommand{\assubtitle}{VERIFY-LAB: Not Yet Run}
\begin{frame}
\frametitle{\astitle}
\begin{block}{\assubtitle}
  \footnotesize
  \begin{itemize}
    \item Prepared without access to Innovus; to check on the first run:
    \begin{itemize}
      \item \cmd{create\_constraint\_mode} / \cmd{create\_analysis\_view}
        on an already initialized design (else: add the scan view to
        \cmd{xh018\_mmmc.tcl} for all steps)
      \item options of \cmd{opt\_design} and \cmd{time\_design} in
        Innovus 26.1
      \item attribute names of the extraction and SI settings
      \item \cmd{add\_metal\_fill} options, density report; X-FAB fill deck?
      \item the slack lines the wrapper reads from the reports
    \end{itemize}
    \item Replace the example output on the Batch slide with a real one
  \end{itemize}
\end{block}
\end{frame}

%%%%%%%%%%%%%%%%%%%%%%%%%%%%%%%%
% Q&A
%%%%%%%%%%%%%%%%%%%%%%%%%%%%%%%%
\section{Q\&A}
\renewcommand{\astitle}{Q\&A}
\begin{frame}
\frametitle{\astitle}
\vfill
\centering
\begin{Huge}\textcolor{red}{Questions?}\end{Huge}
\vfill
\end{frame}

\end{document}
